\documentclass{beamer}
\usetheme{Madrid}
% Use a neutral colortheme so we can fully override colors
\usecolortheme{default}
\usepackage{minted}

% use tikz for generating diagrams
\usepackage{tikz}
\usetikzlibrary{arrows.meta,calc,angles,quotes,positioning}

% for links
\usepackage{hyperref}

% UPenn color palette
\definecolor{PennBlue}{RGB}{1,31,91}   % #011F5B
\definecolor{PennRed}{RGB}{153,0,0}    % #990000

% Apply colors to beamer elements
\setbeamercolor{structure}{fg=PennBlue}
% Make title slide text white for contrast on dark background
\setbeamercolor{title}{fg=white}
\setbeamercolor{subtitle}{fg=black}
\setbeamercolor{author}{fg=black}
\setbeamercolor{date}{fg=black}
% Ensure the title page itself uses our dark background and white foreground
\setbeamercolor{title page}{bg=PennBlue,fg=black}
\setbeamercolor{frametitle}{fg=white,bg=PennBlue}
\setbeamercolor{normal text}{fg=black,bg=white}

% Headline/footline palettes (Madrid uses these)
\setbeamercolor{palette primary}{bg=PennBlue, fg=white}
\setbeamercolor{palette secondary}{bg=PennRed,  fg=white}
\setbeamercolor{palette tertiary}{bg=PennBlue!80!black, fg=white}
\setbeamercolor{palette quaternary}{bg=PennRed!80!black,  fg=white}

% Blocks and emphasis
\setbeamercolor{block title}{fg=white,bg=PennBlue}
\setbeamercolor{block body}{bg=PennBlue!5}
\setbeamercolor{alerted text}{fg=PennRed}
\setbeamercolor{example text}{fg=PennBlue}
\setbeamercolor{itemize item}{fg=PennBlue}
\setbeamercolor{itemize subitem}{fg=PennRed}

% Footer styling colors
\setbeamercolor{author in head/foot}{bg=PennBlue,fg=white}
\setbeamercolor{title in head/foot}{bg=PennRed,fg=white}
\setbeamercolor{date in head/foot}{bg=PennBlue,fg=white}

% Custom footline with course and term centered
\setbeamertemplate{footline}{
  \leavevmode
  \hbox{%%
    \begin{beamercolorbox}[wd=.33\paperwidth,ht=2.5ex,dp=1ex,left]{author in head/foot}
      \hspace{1em}\insertshortauthor
    \end{beamercolorbox}%%
    \begin{beamercolorbox}[wd=.34\paperwidth,ht=2.5ex,dp=1ex,center]{title in head/foot}
      Lecture 12
    \end{beamercolorbox}%%
    \begin{beamercolorbox}[wd=.33\paperwidth,ht=2.5ex,dp=1ex,right]{date in head/foot}
      \insertframenumber{} / \inserttotalframenumber\hspace{1em}
    \end{beamercolorbox}%%
  }
  \vskip0pt
}

% Remove navigation symbols
\setbeamertemplate{navigation symbols}{}


\title{Lecture 12}
\author{ENGR1050 Fall 2025}
\date{\today}

\begin{document}

% Title frame with UPenn logo

\begin{frame}
  \vfill
  \begin{center}

    % Title and course info
    {\Large\textbf{ENGR1050: Lecture 12}}
    \vskip1em
    {\large Exam recap, intro to NumPy, solving linear systems}
    \vskip0.5em
    {\insertdate}
    \vskip2em
    % UPenn logo at bottom
    \includegraphics[height=2.5cm]{upennlogo.png}

  \end{center}
  \vfill
\end{frame}


% Example content frame (optional placeholder)
\begin{frame}{Agenda}
  \begin{itemize}
    \item Go over Exam 1 grades and solution
    \item Some linear algebra review
    \item Introduction to NumPy
  \end{itemize}
\end{frame}

\begin{frame}{Exam 1}

\textbf{Some reminders:}
\begin{itemize}
\item Exam 1 was to assess your grasp of fundamentals before we move into scientific computing
\item Overall, the class did very well. 
\item If you performed less well than you hoped, remember a few things:
    \begin{itemize}
        \item We will continue to use these fundamentals in the course, so you will have more practice
        \item We drop the lowest exam score to make it easier to catch up 
        \item Come to office hours!!!! Many have started coming and made great progress
    \end{itemize}
\end{itemize}

\textbf{New TAs:} We now have two undergraduate TAs who will give you even more chances to get help. They will be in the GM lab 11:30am-3:30pm on Tuesdays. Their names are Kayleen Smith and Annaliese Niebauer and they'll leave a placard on the desk that says \textbf{ENGR1050}.

\end{frame}

\begin{frame}{Exam 1 Grades}
    \begin{columns}
        \begin{column}{0.4\textwidth}
            \textbf{Statistics:} 
            \begin{itemize}
                \item Average: 89
                \item Median: 95
                \item Standard Deviation: 12.3
                \item Highest Score: 100
                \item Lowest Score: 52
            \end{itemize}

            \textbf{Solutions:} \\
            \href{https://colab.research.google.com/github/PIMILab/ENGR1050/blob/main/slides/lecture12/examSolutions.ipynb}{\textcolor{PennBlue}{\textbf{Open in Colab}}}
        \end{column}
        \begin{column}{0.6\textwidth}
            \begin{center}
                \includegraphics[width=1.0\textwidth]{examgrades.png}
            \end{center}
        \end{column}
    \end{columns}
\end{frame}

\begin{frame}{Onward to scientific computing!}
    \begin{itemize}
        \item We will now start learning how to use Python for scientific computing
        \item Today, we will introduce NumPy, the fundamental package for scientific computing in Python
        \item NumPy provides support for arrays (vectors/matrices) and many functions to manipulate them
        \item It is also the basis for many other scientific computing packages in Python, including SciPy, pandas, and matplotlib
    \end{itemize}

    \begin{centering}
    \textit{This unit of the course will be more mathematical - don't get caught off guard! Come to OH early!}
    \end{centering}
\end{frame}

\begin{frame}{Scalars, vectors, and matrices}
    \begin{itemize}
        \item A \textbf{scalar} is a single number, e.g. $x = 5$
        \item A \textbf{vector} is an ordered list of numbers, e.g. $\mathbf{v} = [1, 2, 3]$
        \item A \textbf{matrix} is a 2D array of numbers, e.g. 
        $$\mathbf{M} = \begin{bmatrix}
        1 & 2 & 3 \\
        4 & 5 & 6 \\
        7 & 8 & 9
        \end{bmatrix}$$
        \item Vectors and matrices are fundamental in scientific computing for representing data and performing calculations
    \end{itemize}
\end{frame}

\begin{frame}{Rules to manipulate scalars}
    \begin{itemize}
        \item Addition: $a + b$
        \item Subtraction: $a - b$
        \item Multiplication: $a \times b$ or $a \cdot b$
        \item Division: $\frac{a}{b}$ or $a / b$
        \item Exponentiation: $a^b$ or $a**b$
        \item Order of operations follows standard mathematical rules (PEMDAS/BODMAS)
    \end{itemize}
\end{frame}
\begin{frame}{Rules to manipulate vectors}
    Think of vectors as a list of numbers,
        $$\mathbf{x} = \begin{bmatrix}
    x_1 \\
    x_2 \\
    x_3
    \end{bmatrix}$$ 
    with some additional rules:
    \begin{itemize}
        \item Addition: $\mathbf{u} + \mathbf{v} = [u_1 + v_1, u_2 + v_2, \ldots]$
        \item Subtraction: $\mathbf{u} - \mathbf{v} = [u_1 - v_1, u_2 - v_2, \ldots]$
        \item Scalar multiplication: $c \cdot \mathbf{v} = [c \cdot v_1, c \cdot v_2, \ldots]$
        \item Dot product: $\mathbf{u} \cdot \mathbf{v} = u_1 v_1 + u_2 v_2 + \ldots$
        \item Magnitude (length): $||\mathbf{v}|| = \sqrt{v_1^2 + v_2^2 + \ldots}$
    \end{itemize}
\end{frame}
\begin{frame}{Rules to manipulate matrices}
    Think of matrices as a 2D array of numbers:
    $$\mathbf{A} = \begin{bmatrix}
    a_{11} & a_{12} & a_{13} \\
    a_{21} & a_{22} & a_{23} \\
    a_{31} & a_{32} & a_{33}
    \end{bmatrix}$$
    with some additional rules:
    \begin{itemize}
        \item Addition: $\mathbf{A} + \mathbf{B} = [a_{ij} + b_{ij}]$ (element-wise)
        \item Subtraction: $\mathbf{A} - \mathbf{B} = [a_{ij} - b_{ij}]$ (element-wise)
        \item Scalar multiplication: $c \cdot \mathbf{A} = [c \cdot a_{ij}]$
        \item Matrix multiplication: $\mathbf{A} \mathbf{B}$ (dot product of rows and columns)
        \item Transpose: $\mathbf{A}^T$ (swap rows and columns: $a_{ij} \to a_{ji}$)
    \end{itemize}
    \textit{For some this will be your first encounter with matrices - we will spend time on them in a future lecture. }
\end{frame}


\begin{frame}[fragile]{NumPy scalars/vectors/matrices}
    \begin{itemize}
        \item NumPy provides support for scalars, vectors, and matrices through its \textit{numpy array} object
        \item You can create a NumPy array using the \mintinline{python}{np.array()} function
    \end{itemize}

    \begin{columns}[T,onlytextwidth]
      \begin{column}{0.48\textwidth}
        \begin{minted}[fontsize=\small]{python}
import numpy as np

# Create a 0D array (scalar)
s = np.array(5)

# Create a 1D array (vector)
v = np.array([1, 2, 3])

# Create a 2D array (matrix)
M = np.array([[1, 2, 3],
              [4, 5, 6],
              [7, 8, 9]])
        \end{minted}
      \end{column}
      \begin{column}{0.48\textwidth}
        \begin{minted}[fontsize=\small]{python}
# Element-wise operations
v2 = np.array([4, 5, 6])
v3 = v + 2.3 * v2
# v3 = [1+4*2.3, 
#      2+5*2.3, 
#      3+6*2.3]

        \end{minted}
      \end{column}
    \end{columns}
\end{frame}

\begin{frame}[fragile]{np.array syntax}
\begin{minted}[fontsize=\small]{python}
import numpy as np

# From Python lists
a = np.array([1, 2, 3])              # 1D array
b = np.array([[1, 2], [3, 4]])       # 2D array

# Special arrays
np.zeros((2, 3))     # 2x3 array of zeros
np.ones((3,))        # 1D array of ones
np.eye(3)            # 3x3 identity matrix
np.arange(0, 10, 2)  # [0, 2, 4, 6, 8]
np.linspace(0, 1, 5) # [0., 0.25, 0.5, 0.75, 1.]
\end{minted}
Two ways to construct:
\begin{itemize}
\item From a list for a vector, or list of lists for 2D arrays
\item Using special functions like \mintinline{python}{np.zeros()}, \mintinline{python}{np.ones()}, \mintinline{python}{np.eye()}, \mintinline{python}{np.arange()}, and \mintinline{python}{np.linspace()}
\end{itemize}
\end{frame}


\begin{frame}[fragile]{Accessing np.array properties and elements}
\begin{minted}[fontsize=\small]{python}
a = np.array([1, 2, 3])              # 1D array
b = np.array([[1, 2], [3, 4]])       # 2D array

# Properties of np.arrays
a.shape      # (3,)           → dimensions
a.ndim       # 1              → number of axes
a.size       # 3              → total number of elements
a.dtype      # dtype('int64') → data type

# Indexing/slicing
a[0]         # first element
b[1, 0]      # row 1, column 0
a[1:]        # slice from index 1 to end
b[:, 1]      # all rows, column 1
b[0, :]      # first row

\end{minted}
\begin{itemize}
\item Functions to get shapes of arrays
\item Can index/slice like Python lists, but also with multiple dimensions
\end{itemize}
\end{frame}

\begin{frame}[fragile]{Scalar Multiplication}
We've got a pretty good grip on scalar multiplication:

\begin{minted}[fontsize=\scriptsize]{python}
import numpy as np
A = np.array([2])
B = np.array([3])
C = A * B  # or np.multiply(A, B)
print(C) # [6]
\end{minted}

\end{frame}

\begin{frame}[fragile]{Vector Multiplication}
Comes in two flavors:
 \begin{columns}[T,onlytextwidth]
      \begin{column}{0.50\textwidth}
        \textbf{Scalar multiplication} $\alpha \mathbf{v} = [\alpha v_1, \alpha v_2, \ldots]$
            \vspace{5pt}  
        \includegraphics[width=0.95\textwidth]{scalarproduct.png}  
    \end{column}
      \begin{column}{0.45\textwidth}
        \textbf{Dot multiplication} $\mathbf{u} \cdot \mathbf{v} = u_1 v_1 + u_2 v_2 + \ldots$
        
        \vspace{15pt}  

        \includegraphics[width=0.8\textwidth]{dotproduct.png}
    \end{column}
    \end{columns}
    \centering

\end{frame}

\begin{frame}[fragile]{Vector Multiplication in code}
Comes in two flavors:
 \begin{columns}[T,onlytextwidth]
      \begin{column}{0.50\textwidth}
        \textbf{Scalar multiplication} $\alpha \mathbf{v} = [\alpha v_1, \alpha v_2, \ldots]$
        \begin{minted}[fontsize=\scriptsize]{python}
import numpy as np
alpha = np.array([2])
v = np.array([3, 4, 5])])
out = alpha * v  # or np.multiply(alpha, v)
print(out) # [6, 8, 10]
\end{minted}
      \end{column}
      \begin{column}{0.45\textwidth}
        \textbf{Dot multiplication} $\mathbf{u} \cdot \mathbf{v} = u_1 v_1 + u_2 v_2 + \ldots$
        \begin{minted}[fontsize=\scriptsize]{python}
import numpy as np
u = np.array([2, 3, 4])
v = np.array([5, 6, 7])
out = np.dot(u, v)  # or u @ v
print(out) # 56
\end{minted}
      \end{column}
    \end{columns}
\end{frame}
%============================
\begin{frame}[fragile]{Matrix Multiplication}
\begin{itemize}
    \item Matrix multiplication combines rows of the first matrix with columns of the second.
    \item For $A \in \mathbb{R}^{m\times n}$ and $B \in \mathbb{R}^{n\times p}$, the product $C = AB$ has entries
    \[
        C_{ij} = \sum_{k=1}^{n} A_{ik} B_{kj}.
    \]
    \item Dimensions must satisfy: \textit{(inner dimensions match)} $n=n$.
\end{itemize}
\begin{center}
\begin{minted}[fontsize=\scriptsize]{python}
import numpy as np
A = np.array([[1, 2],
              [3, 4]])
B = np.array([[2, 0],
              [1, 3]])
C = A @ B      # or np.dot(A, B)
print(C)
# [[ 4  6]
#  [10 12]]
\end{minted}
\end{center}
\end{frame}

%============================

\begin{frame}{Example: Writing a System as a Matrix Product}

\begin{block}{Start with two equations}
\[
\begin{aligned}
2x + y &= 5,\\
x - y &= 1.
\end{aligned}
\]
\end{block}

\begin{itemize}
    \item Collect the coefficients of \(x\) and \(y\) into a \textbf{matrix} \(A\):
    \[
    A = \begin{bmatrix} 2 & 1 \\[2pt] 1 & -1 \end{bmatrix},
    \qquad
    \mathbf{x} = \begin{bmatrix} x \\[2pt] y \end{bmatrix},
    \qquad
    \mathbf{b} = \begin{bmatrix} 5 \\[2pt] 1 \end{bmatrix}.
    \]
    \item Multiply \(A\mathbf{x}\) row-by-row:
    \[
    A\mathbf{x} =
    \begin{bmatrix}
        2 & 1\\
        1 & -1
    \end{bmatrix}
    \begin{bmatrix}
        x \\ y
    \end{bmatrix}
    =
    \begin{bmatrix}
        2x + 1y\\
        1x - 1y
    \end{bmatrix}.
    \]
    \item Setting this equal to \(\mathbf{b}\) gives back the original system:
    \[
    \begin{cases}
    2x + y = 5,\\
    x - y = 1.
    \end{cases}
    \]
\end{itemize}
\end{frame}

\begin{frame}[fragile]{Linear Systems as Matrix Equations}

\begin{itemize}
    \item A system of linear equations can be written compactly as
    \[
        A\mathbf{x} = \mathbf{b},
    \]
    where
    \[
    A =
    \begin{bmatrix}
    a_{11} & a_{12} & \cdots & a_{1n}\\
    a_{21} & a_{22} & \cdots & a_{2n}\\
    \vdots & \vdots & \ddots & \vdots\\
    a_{n1} & a_{n2} & \cdots & a_{nn}
    \end{bmatrix},
    \quad
    \mathbf{x} =
    \begin{bmatrix} x_1 \\ x_2 \\ \vdots \\ x_n \end{bmatrix},
    \quad
    \mathbf{b} =
    \begin{bmatrix} b_1 \\ b_2 \\ \vdots \\ b_n \end{bmatrix}.
    \]
    \item Each equation corresponds to \textbf{one row} of \(A\):
    \[
    \begin{aligned}
    a_{11}x_1 + a_{12}x_2 + \cdots + a_{1n}x_n &= b_1,\\
    a_{21}x_1 + a_{22}x_2 + \cdots + a_{2n}x_n &= b_2,\\
    &\;\vdots\\
    a_{n1}x_1 + a_{n2}x_2 + \cdots + a_{nn}x_n &= b_n.
    \end{aligned}
    \]
\end{itemize}
\end{frame}


\begin{frame}[fragile]{Today's exercise: Gaussian elimination}
\begin{itemize}
    \item Goal: Solve a system of linear equations \(A\mathbf{x} = \mathbf{b}\)
    \item Operations allowed (think back to your high school algebra class):
    \begin{enumerate}
        \item Swap two rows.
        \item Multiply a row by a nonzero scalar.
        \item Add a multiple of one row to another.
    \end{enumerate}
\end{itemize}
\end{frame}

%============================
\begin{frame}{Example: A 3$\times$3 Linear System}
Solve
\[
\begin{aligned}
x + y + z &= 6,\\
2x - y + z &= 3,\\
3x + y + 2z &= 11.
\end{aligned}
\]
\textbf{Step 1.} Write as a matrix equation
\[
\left[
\begin{array}{ccc}
1 & 1 & 1 \\
2 & -1 & 1 \\
3 & 1 & 2 
\end{array}
\right] x =
\left[
\begin{array}{c}
6 \\
3 \\
11
\end{array}
\right].
\]
\end{frame}
%============================

\begin{frame}{Example: A 3$\times$3 Linear System}
\[
\left[
\begin{array}{ccc}
1 & 1 & 1 \\
2 & -1 & 1 \\
3 & 1 & 2 
\end{array}
\right] x =
\left[
\begin{array}{c}
6 \\
3 \\
11
\end{array}
\right].
\]

\textbf{Step 2.} Eliminate \(x\) from rows 2 and 3 using row 1.
\[
R_2 \leftarrow R_2 - 2R_1,\quad
R_3 \leftarrow R_3 - 3R_1
\]

\[
\left[
\begin{array}{ccc}
1 & 1 & 1 \\
0 & -3 & -1 \\
0 & -2 & -1 
\end{array}
\right] x =
\left[
\begin{array}{c}
6 \\
-9 \\
-7
\end{array}
\right].
\]
\end{frame}



\begin{frame}{Example: A 3$\times$3 Linear System}
\[
\left[
\begin{array}{ccc}
1 & 1 & 1 \\
0 & -3 & -1 \\
0 & -2 & -1 
\end{array}
\right] x =
\left[
\begin{array}{c}
6 \\
-9 \\
-7
\end{array}
\right].
\]

\textbf{Step 3.} Eliminate \(y\) from row 3 using an integer combination to avoid fractions.
\[
R_3 \leftarrow 3R_3 - 2R_2
\]

\[
\left[
\begin{array}{ccc}
1 & 1 & 1 \\
0 & -3 & -1 \\
0 & 0 & -1 
\end{array}
\right] x =
\left[
\begin{array}{c}
6 \\
-9 \\
-3
\end{array}
\right].
\]
\end{frame}

\begin{frame}{Example: A 3$\times$3 Linear System}
\[
\left[
\begin{array}{ccc}
1 & 1 & 1 \\
0 & -3 & -1 \\
0 & 0 & -1 
\end{array}
\right]
\left[
\begin{array}{c}
x_1  \\
x_2  \\
x_3
\end{array}
\right]
 =
\left[
\begin{array}{c}
6 \\
-9 \\
-3
\end{array}
\right].
\]

\textbf{Step 4.} "Backsolve" by substituting from the bottom up:
\begin{enumerate}
    \item From row 3: 
    \(-x_3 = -3 \Rightarrow x_3 = 3\).
    \item From row 2: 
    \(-3x_2 - x_3 = -9 \Rightarrow -3x_2 - 3 = -9 \Rightarrow x_2 = 2\).
    \item From row 1: 
    \(x_1 + x_2 + x_3 = 6 \Rightarrow x_1 + 2 + 3 = 6 \Rightarrow x_1 = 1\).
\end{enumerate}

\end{frame}

\begin{frame}{Today's exercise}
Repeat the last 4 steps in Python using NumPy.

\href{../notebooks/lec12.ipynb}{Click here to open the Jupyter Notebook for Lecture 12}
\end{frame}
\end{document}